\begin{corollary}[6.1.3]
\EGAReaderTarget{ega4.c04.cor.6-1-3}{IV.2.s6.1.3-fr}
Sean $A$, $B$ dos anillos locales noetherianos, $\mathfrak m$ el ideal maximal
de $A$, $k$ su cuerpo residual y $\varphi:A\to B$ un homomorfismo local.
Supongamos que $\dim(B\otimes_Ak)=0$ \textup{(es decir, (\textbf{0}, 16.2.1), que
$\mathfrak mB$ sea un} ideal de definición \textup{de $B$)}. Entonces se tiene
$\dim(B)\leq\dim(A)$. Si, además, existe un $B$-módulo de tipo finito $N$ que
sea un $A$-módulo plano y cuyo soporte sea igual a $\operatorname{Spec}(B)$
\textup{(lo que ocurre, en particular, si $B$ es un $A$-módulo} plano), se tiene
$\dim(B)=\dim(A)$.
\end{corollary}

La primera afirmación resulta de (5.5.1), y la segunda, de (6.1.2).

\begin{corollary}[6.1.4]
\label{IV.2.s6.1.4-fr}
Sean $X$, $Y$ dos preesquemas localmente noetherianos y $f:X\to Y$ un
morfismo sobreyectivo. Para toda parte cerrada $Z$ de $Y$, se tiene entonces
\begin{equation}
\operatorname{codim}(f^{-1}(Z),X)\leq\operatorname{codim}(Z,Y)
\tag{6.1.4.1}\label{IV.2.s6.1.4.1-fr}
\end{equation}
Además, si $f$ es cuasiplano \textup{(2.3.3)}, los dos miembros de \textup{(6.1.4.1)} son
iguales.
\end{corollary}

En efecto, si $y$ es un punto maximal de $Z$ y $x$ un punto maximal de la
fibra $f^{-1}(y)$, se tiene
$\dim(\mathcal O_{X,x})\leq\dim(\mathcal O_{Y,y})$ en virtud de (5.5.2);
la desigualdad (6.1.4.1) resulta entonces de (5.1.2.1), de (5.1.3.1) y del
hecho de que $f$ es sobreyectivo. Si, además, $f$ es cuasiplano, se sabe
(2.3.4) que toda componente irreducible de $f^{-1}(Z)$ domina una componente
irreducible de $Z$; los puntos maximales de $f^{-1}(Z)$ son, por tanto, los
puntos maximales de las fibras $f^{-1}(y)$, donde $y$ recorre el conjunto de
los puntos maximales de $Z$ $(\mathbf{0}_{\mathrm I}, 2.1.8)$, y en cada uno de estos
puntos maximales $x$ se tiene
$\dim(\mathcal O_x\otimes_{\mathcal O_y}\boldsymbol{k}(y))=0$; como, además,
$\mathcal O_x$ es un $\mathcal O_y$-módulo plano, se tiene
$\dim(\mathcal O_x)=\dim(\mathcal O_y)$ en virtud de (6.1.1), de donde
resulta la igualdad en (6.1.4.1), teniendo en cuenta que $f$ es sobreyectivo.
